\subsection{Inverses}

Let A be a unital Banach algebra, whose identity element is denoted by 1. Recall (TG, IX, prop. 14, p.~40) that the group G of invertible elements of A is an open subset of A, that the topology induced on G by that of A is compatible with the group structure, and that the topological group G is complete.

PROPOSITION 2. --- Let A be a Banach algebra and x an element of A. The series \(\sum_{n=1}^{\infty}\lambda^{n}x^{n}\) considered as a power series in \(\lambda\) , has radius of convergence \(\varrho(x)^{-1}\) . If A is unital and if \(\varrho(x)<1\) , then 1-x is invertible and has inverse \(\sum_{n=0}^{\infty}x^{n}\) .

The series \(\sum_{n=1}^{\infty}\lambda^{n}x^{n}\) has radius of convergence

\[
(\limsup _ {n \to + \infty} \| x ^ {n} \| ^ {1 / n}) ^ {- 1} = \varrho (x) ^ {- 1}
\]

(cf.~VAR, R1, p.~23, 3.1.4). Suppose that A has a unit element. If \(\varrho(x) < 1\) , the series \(\sum_{n=0}^{\infty} x^n\) is therefore absolutely convergent. Since

\[
(1 - x) \left(\sum_ {n = 0} ^ {k} x ^ {n}\right) = \left(\sum_ {n = 0} ^ {k} x ^ {n}\right) (1 - x) = 1 - x ^ {k + 1}
\]

for every integer \(k \geqslant 0\) , the element 1 - x is invertible and its inverse is equal to \(\sum_{n=0}^{\infty} x^{n}\) .

COROLLARY 1. --- If A is a unital Banach algebra, then the group of invertible elements of A contains the open ball with center 1 and radius 1.

This is immediate since \(\|x\| < 1\) implies \(\varrho(x) < 1\) .

COROLLARY 2. --- Let A be a Banach algebra and I a maximal regular left (resp. right) ideal of A. Then I is closed.

Let ( \(\tilde{A}, e\) ) be the unital Banach algebra obtained from A by adjoining an identity element. There exists a maximal left (resp. right) ideal J of \(\tilde{A}\) such that \(J \cap A = I\) (A, VIII, p.~428, prop. 4). Then J is disjoint from the open ball with center e and radius 1 (cor. 1), and hence \(\overline{J} \neq \widetilde{A}\) . Since J is a maximal ideal, this implies that \(\overline{J} = J\) , and consequently \(I = J \cap A = \overline{J} \cap A\) is closed in A.

COROLLARY 3. --- The radical of a Banach algebra is closed.

Indeed, the radical is the intersection of the regular maximal left ideals (A, VIII, p.~430, def. 3).

PROPOSITION 3. --- Let A be a unital Banach algebra.

\begin{enumerate}
\def\labelenumi{\alph{enumi})}
\item
  If \(x \in A\) admits a left inverse (resp. right inverse) y, every element \(x' \in A\) such that \(\|x' - x\| < \|y\|^{-1}\) admits a left inverse (resp. right inverse).
\item
  The set of elements of A that are invertible (resp. left-invertible, resp. right-invertible) is open in A.
\item
  Let \((x_{n})\) be a sequence of elements of A admitting left inverses (resp. right inverses) \(y_{n}\) , such that \((x_{n})\) converges to an element \(x \in A\) and the sequence \((y_{n})\) is bounded. Then \(x\) is left-invertible (resp. right-invertible).
\end{enumerate}

It suffices to treat the case of left inverses; the case of right inverses follows by considering the opposite algebra.

Let \(x, y, x' \in \mathrm{A}\) such that \(yx = 1\) and \(\|x' - x\| < \|y\|^{-1}\) . We have

\[
\left\| 1 - y x ^ {\prime} \right\| = \left\| y x - y x ^ {\prime} \right\| \leqslant \left\| y \right\| \cdot \left\| x - x ^ {\prime} \right\| <   1,
\]

hence \(yx'\) is invertible: there exists \(z \in A\) such that \(z(yx') = 1\) . Thus the element \(x'\) is left-invertible with inverse zy. This proves assertion a), and assertion b) follows immediately.

  Let \((x_{n})\) be a sequence of elements of A admitting left inverses \(y_{n}\) such that \((x_{n})\) converges to an element \(x \in A\) and the sequence \((y_{n})\) is bounded. If \(M \geqslant 1\) is a real number such that \(\|y_{n}\| \leqslant M\) for all \(n \geqslant 1\) , then we have \(\|x_{n} - x\| < M^{-1} \leqslant \|y_{n}\|^{-1}\) for n large enough, and hence x admits a left inverse by a).

DEFINITION 4. --- Let A be a normed algebra, and let x be an element of A. Denote by \(\gamma_x\) and \(\delta_x\) the maps \(y \mapsto xy\) and \(y \mapsto yx\) from A to A. We say that \(x\) is a left (resp. right) topological zero divisor if \(\gamma_x\) (resp. \(\delta_x\) ) is not a homeomorphism from A onto \(\gamma_x(A)\) (resp. onto \(\delta_x(A)\) ).

Remark. --- By TG, IX, p.~36, cor. 2, x is a left (resp. right) topological zero divisor if and only if there exists a sequence \((z_{n})\) in A such that \(\|z_{n}\|=1\) and such that \(xz_{n}\) tends to 0 (resp. that \(z_{n}x\) tends to 0) when \(n\to+\infty\) .

A left (resp. right) topological zero divisor is a left (resp. right) zero divisor. Suppose that A is nonzero and unital. A left (resp. right) topological zero divisor x is not left (resp. right) invertible. Indeed, if for example yx = 1 and if \(xz_{n}\) tends toward 0, then \(z_{n} = y(xz_{n})\) tends to 0 and one cannot have \(\|z_{n}\| = 1\) for all n.

PROPOSITION 4. --- Let A be a unital Banach algebra. Let x be an element of A that is not left-invertible. If there exists a sequence \((x_{n})\) of left-invertible elements of A that converges to x, then x is a right topological zero divisor.

Let \(y_{n}\) be a left inverse of \(x_{n}\) . By prop. 3 (ii), \(\|y_{n}\|\) tends to \(+\infty\) . Let \(z_{n} = \|y_{n}\|^{-1}y_{n}\) . We have \(\|z_{n}\| = 1\) , and \(z_{n}x_{n} = \|y_{n}\|^{-1}\) tends to 0, hence \(z_{n}x = z_{n}x_{n} + z_{n}(x - x_{n})\) tends to 0. One then concludes with the aid of the remark following definition 4.

PROPOSITION 5. --- Let A be a unital Banach algebra and let B be a full subalgebra of A. Then \(\overline{\mathbf{B}}\) is a full subalgebra of A.

Indeed, let \(x\) be an element of \(\overline{\mathrm{B}}\) that is invertible in A, and let \((x_n)\) be a sequence of elements of B tending to \(x\) . Then, for \(n\) sufficiently large, \(x_n\) is invertible in A and \(x_n^{-1}\) tends to \(x^{-1}\) . Since the subalgebra B is full, we have \(x_n^{-1} \in \mathrm{B}\) , whence \(x^{-1} \in \overline{\mathrm{B}}\) .

Let A be a unital Banach algebra and let \((y_{i})_{i\in\mathrm{I}}\) be a family of elements of A. Let B be the full subalgebra of A generated by the elements \(y_{i}\) . Then \(\overline{B}\) is the smallest closed full subalgebra of A containing the \(y_{i}\) . It is called the closed full subalgebra generated by the elements \(y_{i}\) .

